%%
%% ARTICLE.tex — Quantum Encoding Agents
%% ACM sigconf template (acmart class)
%%
%% Compile with:
%%   pdflatex ARTICLE.tex
%%   bibtex   ARTICLE
%%   pdflatex ARTICLE.tex
%%   pdflatex ARTICLE.tex
%%
%% Or on Overleaf: upload this file + references.bib and compile with pdfLaTeX.
%%

\documentclass[sigconf,nonacm]{acmart}

%% ── Packages ───────────────────────────────────────────────────────────────
\usepackage{amsmath}
\usepackage{amssymb}
\usepackage{braket}          % Dirac notation: \ket{}, \bra{}, \braket{}
\usepackage{listings}
\usepackage{xcolor}
\usepackage{graphicx}
\usepackage{booktabs}        % \toprule, \midrule, \bottomrule
\usepackage{hyperref}
\usepackage{microtype}
\usepackage{cleveref}

%% ── Code listing style ─────────────────────────────────────────────────────
\lstset{
  basicstyle=\ttfamily\footnotesize,
  backgroundcolor=\color{gray!8},
  frame=single,
  framerule=0.4pt,
  rulecolor=\color{gray!40},
  breaklines=true,
  captionpos=b,
  keywordstyle=\color{blue!70}\bfseries,
  commentstyle=\color{green!50!black},
  stringstyle=\color{red!60!black},
  showstringspaces=false,
  tabsize=2,
}

%% ── ACM metadata ───────────────────────────────────────────────────────────
%% Remove or fill in when submitting to an actual ACM venue.

\setcopyright{none}
\acmDOI{}
\acmISBN{}
\acmConference[]{}{}{}
\acmYear{2026}
\copyrightyear{2026}

%% ── CCS Concepts ───────────────────────────────────────────────────────────
\begin{CCSXML}
<ccs2012>
 <concept>
  <concept_id>10010147.10010178.10010179</concept_id>
  <concept_desc>Computing methodologies~Quantum computing</concept_desc>
  <concept_significance>500</concept_significance>
 </concept>
 <concept>
  <concept_id>10010147.10010257.10010293.10010294</concept_id>
  <concept_desc>Computing methodologies~Supervised learning</concept_desc>
  <concept_significance>300</concept_significance>
 </concept>
 <concept>
  <concept_id>10010147.10010178.10010219</concept_id>
  <concept_desc>Computing methodologies~Natural language processing</concept_desc>
  <concept_significance>300</concept_significance>
 </concept>
 <concept>
  <concept_id>10002978.10003014.10003017</concept_id>
  <concept_desc>Security and privacy~Access control</concept_desc>
  <concept_significance>100</concept_significance>
 </concept>
</ccs2012>
\end{CCSXML}

\ccsdesc[500]{Computing methodologies~Quantum computing}
\ccsdesc[300]{Computing methodologies~Supervised learning}
\ccsdesc[300]{Computing methodologies~Natural language processing}
\ccsdesc[100]{Security and privacy~Access control}

%% ── Author ─────────────────────────────────────────────────────────────────
\author{Ana Paula Appel}
\email{aappel@redhat.com}
\affiliation{%
  \institution{Red Hat, Inc.}
  \country{USA}
}

%% ── Title ──────────────────────────────────────────────────────────────────
\title{Quantum Encoding Agents: A Natural Language Interface for\\
       Data Embedding Strategy Selection in Quantum Machine Learning}

%% ── Keywords ───────────────────────────────────────────────────────────────
\keywords{Quantum Machine Learning, Quantum Feature Maps, Data Encoding,
          AI Agents, NISQ Hardware, OpenShift, Natural Language Interface}

%% ═══════════════════════════════════════════════════════════════════════════
\begin{document}

\maketitle

%% ── Abstract ───────────────────────────────────────────────────────────────
\begin{abstract}
Selecting an appropriate data encoding strategy is one of the most consequential
and least-tooled decisions in Quantum Machine Learning (QML) pipelines. The choice
of quantum feature map determines the structure of the Hilbert space into which
classical data is embedded, directly shaping the expressibility of quantum kernels,
the trainability of variational circuits, and the feasibility of execution on
near-term hardware. Despite its central role, encoding selection is rarely addressed
systematically: practitioners typically default to a single strategy---most often
angle encoding---without analyzing how the structural characteristics of their data
interact with the properties of each encoding family. This paper presents
\textbf{Quantum Encoding Agents}, an open-source system that reformulates encoding
selection as a natural language interaction. Given a dataset and an optional
description of the QML task, the system analyzes the data profile, applies a
hardware-aware recommendation policy grounded in the critical gate error threshold
$p^{*} \approx 10^{-3}$, generates a complete copyable Qiskit circuit, produces a
natural language justification in the user's language (Portuguese or English), and
computes the quantum kernel matrix with Kernel-Target Alignment scoring. The system
is implemented as a FastAPI microservice, deployed on Red Hat OpenShift with NVIDIA
OpenShell sandbox isolation, and exposed through OpenClaw agents with distinct
epistemic personalities calibrated to different levels of user expertise. Evaluation
confirms that the system correctly selects among seven encoding families---amplitude,
angle, dense angle, IQP, basis, data re-uploading, and custom feature map---under
realistic NISQ hardware constraints, and that KTA scores discriminate between
encodings on benchmark datasets.
\end{abstract}

%% ═══════════════════════════════════════════════════════════════════════════
\section{Introduction}

The practical adoption of Quantum Machine Learning faces a fundamental bottleneck
that precedes model training: the problem of quantum data encoding. Before any
quantum classifier, kernel method, or variational circuit can operate, classical
data must be mapped into quantum states. This mapping---the encoding or feature
map---is not a neutral preprocessing step. It defines the geometry of the state
space, the structure of correlations captured between features, and the depth and
qubit count of the resulting circuit. Different encodings impose radically different
tradeoffs: amplitude encoding compresses $n$ features into $\lceil \log_2 n \rceil$
qubits but requires a state preparation circuit of exponential depth; angle encoding
produces shallow circuits with one qubit per feature but limited expressibility;
IQP encoding~\cite{havlicek2019} captures cross-feature correlations through
diagonal unitaries and has strong theoretical guarantees for quantum kernels, but at
greater circuit depth.

Despite the critical nature of this choice, the QML literature offers little
guidance on encoding selection as a function of data characteristics. Practitioners
encounter a combinatorial decision problem: seven or more encoding families, each
with different qubit efficiency, circuit depth, hardware sensitivity, and affinity
for different QML algorithms, must be matched against a dataset whose relevant
properties---dimensionality, value type, distribution, sign---interact with the
encoding in non-obvious ways. The seminal survey by Sammartino~\cite{sammartino2026},
reviewing 66 papers from 2017 to 2026, identifies this as the most underaddressed
problem in applied QML, formalizes six encoding families, and derives a critical gate
error rate threshold $p^{*} \approx 10^{-3}$ above which deep encodings become
impractical on NISQ hardware.

This paper addresses the encoding selection problem through a different lens: rather
than proposing a new encoding or a new theoretical framework, I build a practical
agent system that makes existing knowledge actionable. I make the following
contributions:

\begin{enumerate}
  \item \textbf{A hardware-aware encoding recommendation engine} that applies the
    $p^{*}$ threshold from~\cite{sammartino2026}, accounting for gate error rate,
    circuit depth budget, qubit count, and chip connectivity topology.

  \item \textbf{Seven implemented encoding strategies} in Qiskit, including
    dense-angle encoding~\cite{sammartino2026} and IQP encoding~\cite{havlicek2019},
    which were identified as underrepresented in applied tooling.

  \item \textbf{A natural language explanation layer} that generates structured
    justifications in Portuguese and English, citing concrete metrics (number of
    qubits, circuit depth, Kernel-Target Alignment) derived from actual simulation
    results.

  \item \textbf{A quantum kernel evaluation endpoint} that computes
    $K[i,j] = |\langle\phi(x_i)|\phi(x_j)\rangle|^2$ for a dataset, returns KTA
    scores, and generates a heatmap visualization---enabling encoding quality
    assessment before any classifier training.

  \item \textbf{A multi-agent deployment architecture} on Red Hat OpenShift with
    NVIDIA OpenShell security sandboxes, exposing three OpenClaw agents with distinct
    communicative personalities for expert and novice audiences.

  \item \textbf{A Kubeflow Pipeline} for Red Hat OpenShift AI that orchestrates the
    full workflow---data analysis, encoding recommendation, comparison, kernel
    evaluation, and MLflow artifact logging---as a reproducible five-stage DAG.
\end{enumerate}

The system is released as open source.\footnote{
  API: \url{https://github.com/anapaulaappel/quantum-encoding-agents};
  Agents: \url{https://github.com/anapaulaappel/openclaw-quantum-agents}.
  All code, policies, agent configurations, and the Kubeflow pipeline were
  developed and are maintained by the author.}

%% ═══════════════════════════════════════════════════════════════════════════
\section{Related Work}

\subsection{Quantum Data Encoding Strategies}

The theoretical landscape of quantum data encoding has been shaped by a sequence of
foundational papers. Schuld and Killoran~\cite{schuld2019} established the
connection between quantum feature maps and kernel methods, showing that quantum
circuits define kernels on data through the inner product of encoded states in
Hilbert space. This framing---encoding as an implicit kernel---is the basis for all
kernel-based QML methods.

Havl\'i\v{c}ek et al.~\cite{havlicek2019} operationalized this idea with a concrete
circuit construction: the ZZ feature map, which applies Hadamard layers followed by
diagonal Pauli rotations encoding both linear ($x_i^2$) and cross-product
($x_i \cdot x_j$) terms. This circuit, equivalent to what we term IQP
(Instantaneous Quantum Polynomial) encoding, was demonstrated on IBM quantum
hardware and showed experimental advantage over classical SVMs on a constructed
classification task.

P\'erez-Salinas et al.~\cite{perezsalinas2020} introduced data re-uploading,
departing from the one-shot encoding paradigm. In their framework, classical data is
inserted into the circuit multiple times across layers, interleaved with trainable
parameters. Data re-uploading has become the de facto encoding for variational
quantum classifiers and quantum neural networks.

LaRose and Coyle~\cite{larose2020} conducted the first systematic comparative study
of encoding strategies across multiple datasets and noise conditions, establishing
that no single encoding dominates and that the optimal choice is dataset-dependent.
Their work motivates the need for an automated selection mechanism---a gap this
system addresses.

The most comprehensive treatment is Sammartino's 2026 survey~\cite{sammartino2026},
which reviews 66 papers from 2017 to 2026, classifies encodings into six families,
and derives practical selection guidelines. Crucially, Sammartino formalizes the
critical gate error threshold: for hardware with gate error rate
$p \geq p^{*} \approx 10^{-3}$, deep encodings (amplitude, IQP, custom feature map)
accumulate noise faster than their expressibility advantage compensates, making
shallow encodings (angle, dense angle, data re-uploading) preferable. Dense-angle
encoding---which packs two features per qubit via
$R_y(x_{2i}) \cdot R_z(x_{2i+1})$---is identified as the most underused encoding
family in applied QML despite its favorable depth-qubit tradeoff.

\subsection{Encoding as an Implicit Kernel}

A result fundamental to the theoretical grounding of this system is Schuld's 2021
theorem~\cite{schuld2021_kernel}: every supervised quantum model that encodes data
into a quantum state $|\phi(x)\rangle$ and measures an observable is mathematically
equivalent to a kernel support vector machine with kernel:
\begin{equation}
  K(x, x') = |\langle\phi(x)|\phi(x')\rangle|^2.
  \label{eq:kernel}
\end{equation}
This equivalence has a decisive practical implication: \emph{``the way that data is
encoded into quantum states is the main ingredient that can potentially set quantum
models apart from classical machine learning models''}~\cite{schuld2021_kernel}.
Furthermore, kernel-based training---which is what QSVM performs directly---is
provably at least as good as variational circuit training for the same encoding. The
choice of encoding is therefore not a preprocessing decision; it is the core
modeling decision, equivalent to choosing a kernel in classical SVM. This system
makes this choice explicit and measurable---the Kernel-Target Alignment (KTA) score
computed by the \texttt{/v1/kernel} endpoint is a direct, label-aware evaluation of
the implicit kernel defined by any encoding on the actual data.

\subsection{Trainability and Barren Plateaus}

The trainability of quantum circuits is a central concern for encoding choices used
in variational algorithms. McClean et al.~\cite{mcclean2018} demonstrated that
randomly initialized quantum circuits exhibit barren plateaus: exponentially
vanishing gradients that make training infeasible beyond moderate qubit counts.
Cerezo et al.~\cite{cerezo2021} surveyed the landscape of variational quantum
algorithms, establishing data re-uploading as the preferred input encoding for QNNs
and VQCs due to its per-layer data injection that maintains gradient flow.

\subsection{Expressibility, Entanglement, and the Encoding Tradeoff}

Sim, Johnson and Aspuru-Guzik~\cite{sim2019} define \emph{expressibility} as the
deviation of a parameterized quantum circuit's output distribution from the Haar
measure. Their key result: shallow circuits have low expressibility and are confined
to structured submanifolds of the state space, which may be insufficient to separate
classes. This is the complementary risk to the barren plateau: too-shallow encodings
underfit in Hilbert space, while too-deep circuits become untrainable due to
vanishing gradients~\cite{mcclean2018}. Sammartino~\cite{sammartino2026} adds the
third axis: above $p^{*} \approx 10^{-3}$, the depth budget is set by hardware
decoherence, not by expressibility requirements.

Larocca et al.~\cite{larocca2023} further establish that over-parameterized quantum
neural networks undergo a phase transition from a barren plateau regime to a
trainable regime as the number of parameters exceeds a critical threshold, directly
linking encoding circuit depth to parameter count and trainability.

\subsection{Quantum Advantage and Scope of NISQ-Era QML}

Huang et al.~\cite{huang2022} proved exponential quantum advantage in a specific
learning task---predicting properties of quantum physical systems from quantum
measurements---demonstrated on 40-qubit superconducting hardware. However, this
advantage is task-specific: it applies when the data itself is quantum in origin,
not when encoding classical tabular data into quantum circuits. Huang, Kueng and
Preskill~\cite{huang2021} establish complementary information-theoretic bounds
showing that classical ML can match quantum ML on average over input distributions,
but quantum advantage is possible worst-case.

For classical datasets on NISQ hardware---the regime this system targets---quantum
advantage over classical methods has not been proven in general. Bowles, Ahmed and
Schuld~\cite{bowles2024} show, in a large-scale study across 12 QML models and 160
datasets, that classical baselines frequently match or outperform quantum classifiers
at current qubit scales, underscoring that encoding evaluation is not optional but
essential. This system is therefore designed as a NISQ-era engineering tool for
informed encoding selection, not as a claim of general quantum supremacy.

\subsection{NISQ vs.\ Fault-Tolerant Quantum Computing}

The \texttt{hardware\_profile} constraints in this system encode NISQ-era realities
that will shift significantly in the fault-tolerant regime. On current NISQ
devices---characterized by gate error rates of $10^{-3}$ to $10^{-2}$, qubit counts
of 10--400, and no error correction---circuit depth is the primary engineering
constraint. In the fault-tolerant era, with logical qubits suppressing error to
arbitrarily low levels, this constraint disappears: amplitude encoding's exponential
qubit compression becomes viable, and the encoding choice is driven purely by
expressibility and task alignment. This design is forward-compatible: passing
\texttt{gate\_error\_rate: 0.0} returns recommendations based on expressibility
alone; passing \texttt{gate\_error\_rate: 5e-3} (IBM Eagle) returns
hardware-adjusted recommendations.

\subsection{AI Agents for Scientific Tooling}

The use of LLM-based agents as interfaces for scientific computing tools is an
emerging paradigm. To my knowledge, this is the first system to deploy
multi-personality LLM agents specifically for quantum encoding recommendation.

The OpenClaw agent framework~\cite{openclaw2026} provides the infrastructure for
agent deployment with persistent memory (via \texttt{MEMORY.md}), user calibration
(via \texttt{BOOTSTRAP.md} and \texttt{USER.md}), and skill encapsulation.

\subsection{Secure Agent Deployment}

NVIDIA OpenShell~\cite{nvidia2026} addresses agent security through out-of-process
policy enforcement: each agent runs in an isolated sandbox governed by a YAML policy
that constrains filesystem access and network endpoints at HTTP method granularity,
enforced by Landlock LSM and seccomp. The policy engine validates rules using the Z3
SMT solver before application. In this deployment, each agent receives a distinct
policy: the expert agent (Circuit) is restricted to the Qiskit API and LLM endpoints
only; the mentor agent (Quanta) additionally has read-only access to arXiv and Qiskit
documentation.

%% ═══════════════════════════════════════════════════════════════════════════
\section{Problem Description}

\subsection{The QML Training Pipeline}

To position the encoding selection problem, I first establish the full QML training
loop in which it occurs. A supervised quantum machine learning pipeline consists of
five stages executed iteratively:

\begin{enumerate}
  \item \textbf{Data Encoding}: $U(x)|0\rangle^{\otimes n} \rightarrow |\phi(x)\rangle$
        \hfill \emph{(this paper)}
  \item \textbf{Parameterized Ansatz}: $V(\theta)|\phi(x)\rangle \rightarrow |\psi(x,\theta)\rangle$
  \item \textbf{Measurement}: $\langle\psi(x,\theta)|M|\psi(x,\theta)\rangle \rightarrow \hat{y}$
  \item \textbf{Classical Loss}: $\mathcal{L}(\hat{y}, y) \rightarrow \mathbb{R}$
  \item \textbf{Classical Optimizer}: $\theta \leftarrow \theta - \eta\nabla\mathcal{L}$
        (then repeat)
\end{enumerate}

The encoding layer (stage 1) is fixed for the duration of training---it is a
hyperparameter of the model, not a trainable parameter. A poorly chosen encoding
defines a kernel that cannot linearly separate the classes in Hilbert space, no
matter how many training iterations the optimizer runs. This system operates at
stage 1, with KTA providing a measurement of whether the chosen encoding creates a
Hilbert space geometry useful for the loss at stage 4.

For quantum kernel methods (QSVM), stages 2--5 are replaced by classical SVM
training on the kernel matrix $K[i,j] = |\langle\phi(x_i)|\phi(x_j)\rangle|^2$---
which is exactly what the \texttt{/v1/kernel} endpoint computes. In this case, the
encoding is the \emph{entire} model specification, reinforcing the centrality of
stage 1.

\subsection{The Encoding Selection Problem}

Let $\mathcal{D} = \{x_i\}_{i=1}^N \subset \mathbb{R}^d$ be a classical dataset
with $N$ samples and $d$ features. A quantum encoding is a parameterized map
$\phi: \mathbb{R}^d \rightarrow \mathcal{H}$ from the feature space to a Hilbert
space $\mathcal{H}$ of dimension $2^n$, implemented as a quantum circuit $U(x)$
acting on $n$ qubits such that $|\phi(x)\rangle = U(x)|0\rangle^{\otimes n}$.

The encoding selection problem is: given $\mathcal{D}$ and an optional QML task
specification $\mathcal{T}$ (classification, kernel method, variational circuit,
etc.)\ and hardware constraints $\mathcal{H}\!w$ (gate error rate, qubit count,
connectivity topology), select an encoding $\phi^{*}$ that maximizes downstream task
performance subject to hardware feasibility.

This problem is ill-posed for two reasons. First, the downstream task performance
depends on the full learning pipeline, not just the encoding; it cannot be evaluated
without training a model. Second, the interaction between data characteristics and
encoding properties is non-linear and high-dimensional.

\subsection{Current Practice and Its Limitations}

In practice, encoding selection is rarely principled. A survey of QML
implementations in the literature reveals that the majority use angle encoding as a
default, regardless of dataset structure. However, this choice ignores:

\begin{itemize}
  \item \textbf{Dimensionality}: for $d > 10$ features, angle encoding requires $d$
    qubits. Dense-angle encoding achieves the same depth with $\lceil d/2 \rceil$
    qubits.

  \item \textbf{Task alignment}: for QSVM, encodings without entanglement define
    kernels with limited expressibility. IQP and custom feature map encodings define
    richer kernels.

  \item \textbf{Hardware noise}: above $p^{*} \approx 10^{-3}$, amplitude encoding's
    deep state preparation circuit accumulates noise faster than its qubit compression
    advantage is worth. On IBM Eagle hardware ($p \approx 5 \times 10^{-3}$), the
    effective circuit depth after error accumulation can reduce the signal-to-noise
    ratio below classification threshold.

  \item \textbf{Data type}: basis encoding, the natural choice for binary or
    categorical data, is systematically ignored when practitioners default to angle
    encoding for all data types.
\end{itemize}

\subsection{The Accessibility Gap}

Beyond the technical selection problem, there is an accessibility gap: the knowledge
required to make principled encoding choices---distributed across a dozen papers
published between 2018 and 2026, requiring familiarity with quantum information
theory, circuit complexity, and NISQ hardware characteristics---is not accessible to
the ML practitioner approaching QML for the first time.

This gap manifests in two distinct user profiles. The \textbf{QML practitioner}
(expert) needs dense, metric-first information: encoding name, circuit depth, qubit
count, KTA score, and a reference to the relevant paper. The \textbf{ML practitioner
approaching QML} (learner) needs progressive disclosure: physical intuition before
mathematics, the Bloch sphere before Dirac notation, a concrete analogy before a
formal definition. Current tools address neither profile. The system I present
bridges this gap through specialized agents calibrated to each profile.

%% ═══════════════════════════════════════════════════════════════════════════
\section{Quantum Encoding Agents: System Design}

\subsection{Architecture Overview}

The system consists of three loosely coupled layers: (1) an OpenClaw agent gateway
with LLM backend (Qwen2.5-Coder-14B via Ollama or Llama-3.2-3B via vLLM on
OpenShift), (2) the \texttt{llama-qiskit-agents} FastAPI microservice exposing the
quantum pipeline, and (3) the NVIDIA OpenShell security layer enforcing per-agent
sandbox policies. The microservice is stateless and can be called directly via HTTP
independently of the agent framework.

\subsection{Data Profile Analysis}

The first processing step is structured data profiling. Given input in any of three
forms---a numpy array, a CSV file, or a natural language description---the system
extracts a \texttt{DataProfile} dataclass with seven fields:
\texttt{n\_samples}, \texttt{n\_features}, \texttt{is\_binary},
\texttt{is\_categorical}, \texttt{is\_continuous}, \texttt{has\_negative}, and
\texttt{description}. For text descriptions, keyword detection in both Portuguese
and English maps phrases to profile flags.

\subsection{The Seven Encoding Strategies}

The system implements seven encoding families as Qiskit \texttt{QuantumCircuit}
objects. Table~\ref{tab:encodings} summarizes their properties.

\begin{table}[h]
\caption{Quantum encoding strategies implemented in the system.}
\label{tab:encodings}
\begin{tabular}{lllp{3.2cm}l}
\toprule
\textbf{Encoding} & \textbf{Qubits} & \textbf{Depth} &
\textbf{Recommended for} & \textbf{Key gates} \\
\midrule
Amplitude       & $\lceil\log_2 n\rceil$ & $O(2^n)$  & Large vectors, qubit-constrained & \texttt{StatePreparation} \\
Angle           & $d$                    & $1$        & $d\!\leq\!4$, continuous         & $R_y(x_i)$ \\
Dense Angle~\cite{sammartino2026} & $\lceil d/2\rceil$ & $2$ & $5\!\leq\!d\!\leq\!12$, contin.\ & $R_y R_z$ per qubit \\
IQP~\cite{havlicek2019} & $d$            & ${\sim}3d$ & $8\!\leq\!d\!\leq\!16$, kernels  & $H{+}R_z(x_i^2){+}R_{zz}$ \\
Basis           & $d$                    & ${\leq}d$  & Binary / categorical             & $X$ per 1-bit \\
Data Re-upl.~\cite{perezsalinas2020} & $d$ & ${\sim}4d$ & VQC, QNN              & $R_y{\times}L + CX$ chain \\
Custom FM       & $d$                    & ${\sim}3d$ & QSVM, arbitrary kernels          & $H{+}R_z{+}R_y{+}CZ$ \\
\bottomrule
\end{tabular}
\end{table}

\textbf{Dense-angle encoding}~\cite{sammartino2026} achieves half the qubit count of
standard angle encoding at depth 2 by packing two features per qubit:
\begin{equation}
  U_{\mathrm{DA}}(x)\,|0\rangle^{\otimes \lceil d/2 \rceil}
  = \bigotimes_{i=0}^{\lceil d/2 \rceil - 1}
    R_z(x_{2i+1})\,R_y(x_{2i})\,|0\rangle_i.
\end{equation}

\textbf{IQP encoding}~\cite{havlicek2019} is implemented without a
\texttt{qiskit-machine-learning} dependency through manual decomposition of
$R_{zz}(\theta) = CX \cdot R_z(\theta) \cdot CX$:
\begin{equation}
  U_{\mathrm{IQP}}(x)
  = H^{\otimes d}
    \cdot \prod_i R_z(x_i^2)
    \cdot \prod_{i<j} R_{zz}(x_i x_j)
    \cdot H^{\otimes d}.
\end{equation}
Only adjacent pairs $(i, i{+}1)$ are connected, keeping circuit depth polynomial in
$d$ and compatible with linear and heavy-hex topologies.

\subsection{Hardware-Aware Recommendation Policy}

The recommendation pipeline has two stages: data-driven base selection and
hardware-constrained refinement.

\textbf{Base selection} applies a priority decision tree over the
\texttt{DataProfile}: (1) binary, $d \leq 16$ $\rightarrow$ \texttt{basis};
(2) continuous, $d \leq 4$ $\rightarrow$ \texttt{angle};
(3) continuous, $4 < d \leq 12$, no negatives $\rightarrow$ \texttt{dense\_angle};
(4) single sample, $d \geq 4$ $\rightarrow$ \texttt{amplitude};
(5) continuous, $8 < d \leq 16$ $\rightarrow$ \texttt{iqp};
(6) continuous, $d > 16$ $\rightarrow$ \texttt{custom\_feature\_map};
(7) default $\rightarrow$ \texttt{angle}.

\textbf{Task refinement} overrides based on the QML task: kernel methods and QSVM
elevate to \texttt{custom\_feature\_map}; variational methods prefer
\texttt{data\_reuploading}; binary classification with binary data uses
\texttt{basis}.

\textbf{Hardware refinement} applies the $p^{*}$ threshold from
Sammartino~\cite{sammartino2026}. For any \texttt{HardwareProfile} with
\texttt{gate\_error\_rate} $\geq 10^{-3}$, deep encodings
(\texttt{amplitude}, \texttt{custom\_feature\_map}, \texttt{iqp}) are replaced by
\texttt{dense\_angle} ($d > 4$) or \texttt{angle} ($d \leq 4$). Additionally,
\texttt{connectivity: heavy-hex} or \texttt{linear} triggers a SWAP overhead
warning for \texttt{custom\_feature\_map}, and \texttt{max\_depth\_budget} triggers
a depth feasibility warning.

\subsection{Natural Language Explanation Generation}

The explanation layer generates a four-paragraph structured narrative in the
detected language (Portuguese or English), each paragraph citing concrete values
from the \texttt{DataProfile} and \texttt{SimulationResult}: (1) data description,
(2) encoding justification with geometric intuition, (3) circuit metrics with NISQ
compatibility assessment, and (4) comparative analysis of alternatives with concrete
depth and qubit differences. Language detection uses vocabulary overlap between the
input text and Portuguese/English marker sets.

\subsection{Quantum Kernel Evaluation}

For datasets where classification labels are available, the system computes the
quantum kernel matrix $K \in [0,1]^{N \times N}$:
\begin{equation}
  K[i,j] = |\langle\phi(x_i)|\phi(x_j)\rangle|^2.
\end{equation}
Statevectors are computed using the \texttt{StatevectorSimulator} (exact, without
measurement shots) and inner products are computed classically. The
Kernel-Target Alignment is computed as:
\begin{equation}
  \mathrm{KTA}(K, y)
  = \frac{\langle K,\, yy^{\top} \rangle_F}
         {\|K\|_F \cdot \|yy^{\top}\|_F},
\end{equation}
where $y \in \{-1, +1\}^N$ is the label vector. KTA provides a label-aware measure
of how well the encoding separates classes in Hilbert space, independently of any
trained classifier. The result is returned as a JSON matrix, a heatmap visualization
(viridis colormap, per-class tick coloring), and a natural language caption.

In a preliminary evaluation on a synthetic two-class dataset ($N=4$, $d=2$), angle
encoding achieved KTA $= 0.701$ while custom feature map achieved KTA $= 0.114$,
demonstrating that richer encodings do not necessarily produce better-aligned kernels
and that empirical evaluation via KTA is essential before committing to a training
pipeline.

\subsection{Bloch Sphere Visualization}

For encodings with $n \leq 6$ qubits, the system optionally captures the statevector
before measurement and renders the Bloch sphere representation of each qubit's
marginal state using \texttt{plot\_bloch\_multivector}. The resulting PNG
(base64-encoded) accompanies the recommendation when \texttt{include\_bloch: true}
is set in the request, providing geometric intuition for what the encoding does to
each data point.

\subsection{Agent Personalities and the Role of SOUL.md}

Three OpenClaw agents expose the recommendation system through distinct communicative
registers. \textbf{QiskitAgent} (\texttt{quantum\_encoding}) is the general-purpose
agent. \textbf{Circuit} (\texttt{qiskit\_expert}) is calibrated for QML
practitioners: its \texttt{SOUL.md} instructs it to presuppose familiarity with key
concepts, respond in three dense blocks (encoding name, metrics, rationale), never
explain superposition, and present trade-offs as two objective sides without
choosing. \textbf{Quanta} (\texttt{qiskit\_mentor}) is calibrated for learners: its
\texttt{SOUL.md} prohibits the Schr\"odinger's cat analogy (which lies) and mandates
physical intuition before mathematics---the Bloch sphere before Dirac notation. Both
agents perform a calibration ritual on first session (\texttt{BOOTSTRAP.md}) to
populate \texttt{USER.md} with a persistent user profile.

\subsection{Skill Architecture}

Two shared skills encapsulate reusable behavior across all agents.
\texttt{\$qiskit-api} documents the \texttt{/v1/recommend/explain} endpoint as a
\texttt{SKILL.md} file, eliminating documentation duplication and ensuring
consistent behavior when the API evolves. \texttt{\$circuit-review} provides a
four-step protocol for reviewing an existing circuit pasted by the user: identify
the encoding, extract metrics, call \texttt{/v1/compare} for benchmark, and issue
one of three verdicts (keep / optimize / replace) with a technical rationale.

\subsection{Security Architecture (OpenShell)}

Each agent runs inside a dedicated NVIDIA OpenShell sandbox governed by a YAML
policy validated by the Z3 SMT solver. The policies enforce least-privilege at HTTP
method granularity: Circuit (expert) has access limited to
\texttt{qiskit-api:8080} (read-write) and the LLM endpoint (POST only); Quanta
(mentor) additionally has read-only access to \texttt{arxiv.org} and
\texttt{docs.quantum.ibm.com}; QiskitAgent (general) includes write access to
\texttt{/tmp} for CSV uploads. Binaries are specified as objects with \texttt{path}
keys, enabling per-binary policy enforcement through Landlock LSM.

\subsection{Kubeflow Pipeline for RHOAI}

The full workflow is codified as a Kubeflow Pipelines v2 DAG deployable on Red Hat
OpenShift AI~\cite{rhoai2026}: (1) \texttt{analyze\_data} via
\texttt{POST /v1/analyze}, (2) \texttt{recommend\_encoding} via
\texttt{POST /v1/recommend/explain}, (3) \texttt{compare\_encodings} via
\texttt{POST /v1/compare}, (4) \texttt{compute\_kernel} via \texttt{POST /v1/kernel}
(when labels are provided), and (5) \texttt{log\_to\_mlflow}~\cite{zaharia2018}
(parameters, metrics, heatmap, Bloch sphere visualization, generated Qiskit code).
The pipeline accepts hardware profile parameters that propagate to the
recommendation engine, enabling reproducible NISQ-aware experiments.

%% ═══════════════════════════════════════════════════════════════════════════
\section{Discussion}

\subsection{The Role of KTA in Practice}

My preliminary results suggest that KTA is a more reliable pre-training signal than
circuit expressibility alone. Angle encoding (KTA $= 0.701$) outperformed the
custom feature map (KTA $= 0.114$) on the test dataset despite the latter's greater
theoretical richness. This confirms LaRose and Coyle's finding~\cite{larose2020}
that no encoding universally dominates, and motivates the design approach: rather
than recommending based on theoretical properties alone, I compute KTA on the actual
data and present the score alongside the recommendation.

\subsection{Limits of Text-Based Encoding Selection}

The natural language input path---inferring data profile from a description---has
inherent limitations. Keyword detection can misclassify ambiguous descriptions. For
production use, numerical data or CSV upload should be preferred; the text path is
primarily useful for initial exploration and for the conversational interface.

\subsection{Agent Personality as an Engineering Artifact}

The \texttt{SOUL.md}-based personality specification is a form of behavioral
engineering that sits between prompt engineering and agent architecture. Unlike a
system prompt, \texttt{SOUL.md} is persistent (injected every session), versioned
in git, and testable: a response from Circuit that begins with ``Great question!''
violates the soul specification and can be detected in evaluation. This framing---
soul as contract---enables more systematic quality assurance of agent behavior than
ad-hoc prompting.

\subsection{OpenShell Policy Design as Security-by-Data-Profile}

The assignment of different network permissions to different agents based on their
communicative role is a novel security design pattern: \textbf{policy as a function
of agent persona}. Quanta requires external internet access (arXiv, Qiskit docs)
because its role is pedagogical and citation-heavy. Circuit does not, because
experts are expected to know the literature. This reflects the insight that the
minimum necessary privilege of an AI agent is determined not just by its technical
function but by its communicative purpose.

%% ═══════════════════════════════════════════════════════════════════════════
\section{Conclusion}

I presented Quantum Encoding Agents, a system that addresses the encoding selection
problem in Quantum Machine Learning through a combination of hardware-aware
recommendation logic, natural language generation, quantum kernel evaluation, and
specialized AI agents calibrated to different user expertise levels. The system is
grounded in the theoretical framework of Sammartino~\cite{sammartino2026} and
Havl\'i\v{c}ek et al.~\cite{havlicek2019}, implements seven encoding families
including the underused dense-angle encoding, and provides the first open-source
implementation of KTA-based encoding quality assessment via a REST API.

The multi-agent architecture demonstrates that AI agents deployed for scientific
tools can and should have distinct communicative personalities calibrated to their
user populations---not as a cosmetic feature, but as a substantive design choice
that determines what information to surface, in what order, and with what degree of
assumed prior knowledge. The security architecture extends this idea: agent policy
is a function of agent role, not just of agent function.

Future work includes: (i) integration with IBM Quantum real hardware via
\texttt{QiskitRuntimeService} for empirical NISQ validation; (ii) the Q-Tucker
shallow state preparation~\cite{qtuckerarxiv} as an alternative amplitude encoding
with $O(n)$ depth; (iii) the Quantum Spectral Model~\cite{qsmarxiv} as an adaptive
encoding whose frequency spectrum is conditioned on the input; and (iv) a
KTA-guided encoding search that evaluates multiple candidates and returns the one
with the highest alignment score.

%% ═══════════════════════════════════════════════════════════════════════════
\bibliographystyle{ACM-Reference-Format}
\bibliography{references}

\end{document}
